% Chapter 2
\section{Trabalhos Relacionados}
\label{sec:trabalhos_correlatos}

O desenvolvimento de soluções automatizadas para o agendamento de aulas e exames tem sido impulsionado pela aplicação de heurísticas e meta-heurísticas. Diversos estudos indicam que essas abordagens são adequadas para lidar com o amplo espaço de busca característico dos problemas de timetabling, especialmente em cenários nos quais a enumeração exaustiva se torna inviável.

Bashab et al. (2023) apresentam uma revisão abrangente sobre técnicas de otimização aplicadas à calendarização universitária, contemplando conceitos, metodologias, restrições, benchmarks e questões em aberto. O trabalho destaca que o university timetabling é tratado na literatura como um problema de otimização combinatória e de natureza não polinomial, no qual diferentes estratégias computacionais têm sido empregadas para buscar soluções de boa qualidade. Entre essas estratégias, os autores observam o uso recorrente de meta-heurísticas, métodos multiobjetivo e abordagens híbridas, especialmente pela capacidade dessas técnicas de tratar múltiplas restrições e critérios de avaliação em instâncias complexas.

\citeonline{awad2022} propuseram uma abordagem baseada em Busca Tabu Adaptativa para problemas de course timetabling em larga escala. A principal contribuição do estudo está no ajuste dinâmico da lista tabu, utilizado para melhorar o processo de busca e evitar estagnação em regiões pouco promissoras do espaço de soluções. A validação foi realizada com onze instâncias de benchmark originalmente apresentadas por Socha et al., compostas por problemas pequenos, médios e grandes. Os resultados indicaram que a abordagem foi capaz de produzir grades competitivas em comparação com métodos anteriores, reforçando a relevância da adaptação de parâmetros em meta-heurísticas aplicadas à calendarização acadêmica.

\citeonline{elen2022}, por sua vez, demonstrou a aplicação prática de Algoritmos Genéticos no problema de agendamento de exames. O estudo desenvolveu um método baseado em AG para o Vocational School of T.O.B.B. Technical Sciences, da Karabuk University, utilizando dados reais do período letivo de 2018-2019. A solução foi testada com 421 disciplinas, 54 docentes, 27 programas acadêmicos e 21 salas, obtendo horários de exame sem sobreposição em tempo médio de aproximadamente quatro minutos e em menos de 2000 iterações. Esse resultado evidencia a viabilidade do uso de Algoritmos Genéticos em sistemas de apoio à decisão voltados ao planejamento acadêmico.

Em conjunto, esses trabalhos demonstram que meta-heurísticas como Algoritmos Genéticos, Busca Tabu e abordagens híbridas constituem alternativas relevantes para problemas de calendarização acadêmica. Entretanto, observa-se que muitos estudos concentram-se na avaliação algorítmica isolada, sem necessariamente integrar a solução a um sistema operacional completo ou comparar diferentes paradigmas de implementação sob uma mesma arquitetura. Nesse ponto, o IGRIS se diferencia ao combinar a geração automatizada de grades com módulos de gestão acadêmica, interface de operação e benchmark controlado entre implementações desenvolvidas em Java, Python e Go.
